\chapter{Úvod}
\label{sec:introduction}

Počítačové sítě jsou dnes základní infrastrukturou, na jejíž spolehlivosti závisí provoz institucí i výrobních podniků. S rostoucím počtem připojených zařízení a s postupným zaváděním protokolu IPv6 vedle zavedeného protokolu IPv4 roste i složitost úloh, které musí správce sítě řešit přímo na místě. Otázky bývají přitom velmi konkrétní: zda je zásuvka v místnosti skutečně propojena s rozvaděčem, jakou propustnost daná linka nabízí, která zařízení se v segmentu nacházejí a zda na nich neběží služby, které by běžet neměly.

Odpověď na tyto otázky vyžaduje měření na místě. Právě zde naráží běžně dostupná řešení na svá omezení. Komerční příruční testery jsou ergonomické, ale uzavřené, takže rozsah jejich funkcí je pevně dán výrobcem a vlastní testovací postup do nich doplnit nelze. Notebook s otevřenými nástroji je naopak plně flexibilní, avšak v rozvaděči nebo pod stolem se s ním pracuje obtížně a před každým měřením vyžaduje ruční nastavení. Pokusí-li se navíc technik připojit do téhož segmentu dvěma rozhraními, aby mohl současně naslouchat provozu a generovat zátěž, narazí na konflikt ve směrovacích tabulkách operačního systému, který výsledky měření nenápadně zkreslí.

Cílem této bakalářské práce je navrhnout a realizovat přenosné zařízení pro testování a monitorování počítačových sítí, které spojuje otevřenost softwarového řešení s ergonomií samostatného přístroje. Prototyp je postaven na jednodeskovém počítači Raspberry Pi, je vybaven dotykovým displejem a vlastní řídicí aplikací, takže jej lze ovládat bez připojeného počítače a bez fyzické klávesnice. Umožňuje pasivní analýzu provozu, aktivní měření propustnosti, mapování segmentu i bezpečnostní průzkum, a to shodně pro protokoly IPv4 a IPv6.

Klíčovým technickým prvkem návrhu je logické oddělení dvou testovacích rozhraní pomocí technologie Linux Network Namespaces. Toto oddělení zajišťuje, že generovaná zátěž neovlivňuje statistiky zachycené naslouchajícím rozhraním, a je vynuceno přímo jádrem operačního systému, nikoli pouze konfigurační kázní obsluhy. Druhou oblastí, které je věnována zvýšená pozornost, je automatická konfigurace adres protokolu IPv6, neboť právě na ní závisí, zda přístroj dokáže navázat spojení se stanicí, na níž není nic předem nastaveno.

Text práce je rozdělen do čtyř kapitol. Po tomto úvodu následuje kapitola \ref{sec:teorie}, která shrnuje teoretická východiska: rozdíly mezi protokoly síťové vrstvy, adresaci a automatickou konfiguraci protokolu IPv6, nástroje pro monitorování a audit provozu a architekturu jmenných prostorů. Její závěrečná část porovnává existující řešení a formuluje požadavky na navrhované zařízení. Kapitola \ref{sec:realizace} popisuje návrh přístroje po hardwarové i softwarové stránce, realizaci generátoru a analyzátoru provozu a zhodnocení výsledků ověřených laboratorním měřením. Závěr práce shrnuje dosažené výsledky, vymezuje vlastní přínos autora a naznačuje směry dalšího vývoje prototypu.